\documentclass[11pt,a4paper]{article}
\usepackage[utf8]{inputenc}
\usepackage[spanish]{babel}
\usepackage{amsmath,amsfonts,amssymb}
\usepackage{booktabs}
\usepackage{geometry}
\geometry{margin=2.5cm}

\title{Ecuaciones del modelo Dynamic Energy Budget (DEB) \\ aplicado a larvas de insectos}
\author{Resumen extraído de la biblioteca Zotero local}
\date{\today}

\begin{document}
\maketitle

\section*{Nota previa sobre los PDFs consultados}

Se revisaron los archivos PDF de la carpeta \texttt{/home/jjlealg/Zotero}. Los únicos documentos que abordan explícitamente un modelo tipo DEB (Dynamic Energy Budget) aplicado a larvas/insectos son dos artículos sobre \textit{Hermetia illucens} (mosca soldado negra):

\begin{itemize}
    \item Padmanabha et al.\ (2020). \textit{A comprehensive dynamic growth and development model of Hermetia illucens larvae}. PLOS ONE.
    \item Eriksen (2022). \textit{Dynamic modelling of feed assimilation, growth, lipid accumulation, and CO\textsubscript{2} production in black soldier fly larvae}. PLOS ONE (archivo \texttt{journal.pone.0276605}).
    \item Eriksen (2024). \textit{Metabolic performance and feed efficiency of black soldier fly larvae}.
\end{itemize}

Ninguno de los PDFs encontrados implementa el DEB ``canónico'' de Kooijman con las variables exactas $E$ (reserva) y $V$ (volumen estructural). Ambos son \textbf{modelos inspirados en DEB} que particionan la biomasa en compartimentos análogos. A continuación se presenta: (i) el marco teórico estándar DEB con $E$ y $V$, y (ii) las ecuaciones extraídas de los PDFs, destacando la correspondencia conceptual.

% --------------------------------------------------------------------
\section{1. Marco teórico estándar DEB (Kooijman, 2010)}
% --------------------------------------------------------------------

En el DEB clásico, el organismo se describe mediante dos variables de estado: la energía en reserva $E(t)$ y el volumen estructural $V(t)$. Sus ecuaciones diferenciales principales son:

\begin{align}
    \frac{dE}{dt} &= p_A - p_C \label{eq:deb_E} \\[6pt]
    \frac{dV}{dt} &= \frac{\kappa p_C - p_M}{[E_G]} \label{eq:deb_V}
\end{align}

donde la tasa de utilización $p_C$ se reparte en mantenimiento, crecimiento y maduración/reproducción:

\begin{equation}
    p_C = [p_M]V + [E_G]\frac{dV}{dt} + \frac{p_R}{\kappa}
    \label{eq:deb_pc}
\end{equation}

Frecuentemente se escribe también la dinámica de la densidad de reserva $e=E/V$:

\begin{equation}
    \frac{de}{dt} = \frac{\{p_{Am}\}f - [p_M]e}{e + [E_G]} - e\,\frac{\kappa p_C - p_M}{V[E_G]}
    \label{eq:deb_e}
\end{equation}

aunque la forma más usada en aplicaciones empíricas es el sistema $(E,V)$ dado por (\ref{eq:deb_E})--(\ref{eq:deb_V}).

% --------------------------------------------------------------------
\section{2. Modelo DEB-inspired de Padmanabha et al.\ (2020) para \textit{H.~illucens}}
% --------------------------------------------------------------------

Este trabajo no usa $E$ y $V$ directamente, sino que modela la \textbf{masa seca por larva} $B_{dry}(t)$ [g] y la \textbf{suma de desarrollo} $TS(t)$ [h]. Conceptualmente:

\begin{itemize}
    \item $B_{dry}$ agrupa la masa estructural y las reservas (análogo combinado de $E+V$).
    \item $B_{eff}$ = ``asimilados no estructurales'' (análogo de la reserva $E$).
    \item $B_{str}$ = ``masa estructural del cuerpo'' (análogo del volumen estructural $V$).
\end{itemize}

La ecuación de balance de masa general es:

\begin{equation}
    \frac{dB_{dry}}{dt} = \dot{B}_{ing}
    - \dot{B}_{excr}
    - \dot{B}_{assim}
    - \dot{B}_{maint}
    - \dot{B}_{mat}
    - \dot{B}_{metab}
    \label{eq:pad_balance}
\end{equation}

Definiendo los asimilados efectivos disponibles:

\begin{equation}
    \dot{B}_{eff} = (1 - k_{a,excr} - k_{a,assim})\,\dot{B}_{ing}
    = \eta_{inges}\,k_{inges}\,B_{dry}
    \label{eq:pad_beff}
\end{equation}

y agrupando mantenimiento y madurez en un solo flujo, el modelo dinámico queda:

\begin{align}
    \frac{dB_{dry}}{dt} &=
    \eta_{inges}\,r_{assim}(T, TS, B_{feed}, W_{med}, A_{air})\,k_{inges}\,B_{dry}
    - r_{mat}(T, TS, B_{feed}, A_{air})\,k_{maint}\,B_{dry}
    \label{eq:pad_bdry} \\[6pt]
    \frac{dTS}{dt} &= r_{dev}(T_{med}, B_{feed}, W_{med}, A_{air})\,k_{devts}
    \label{eq:pad_ts}
\end{align}

Las funciones de regulación $r_{assim}$, $r_{mat}$ y $r_{dev}$ incorporan los efectos de temperatura ($r_T$), disponibilidad de alimento ($r_F$/$r_{Fgrw}$), humedad ($r_W$) y aireación ($r_A$). Por ejemplo:

\begin{align}
    r_{assim} &= r_{Bassim}(TS)\,
    \frac{r_T(T_{med})}{k_{rmax,T}}\,
    \frac{r_{Fgrw}(B_{feed})}{k_{rmax,gm}}\,
    \frac{r_W(W_{med})}{k_{rmax,W}}\,
    \frac{r_A(A_{air})}{k_{rmax,A}}
    \label{eq:pad_rassim} \\[6pt]
    r_{mat} &= r_{Bmat}(TS)\,
    \frac{r_T(T_{med})}{k_{rmax,T}}\,
    \frac{r_{Fgrw}(B_{feed})}{k_{rmax,gm}}\,
    \frac{r_A(A_{air})}{k_{rmax,A}}
    \label{eq:pad_rmat}
\end{align}

La función de etapa de alimentación $r_{Bassim}(TS)$ decrece logísticamente desde el máximo hasta cero cuando el desarrollo supera $k_{TS2}$; $r_{Bmat}(TS)$ representa la asignación a madurez y es la que produce la caída de masa observada en prepupas.

% --------------------------------------------------------------------
\section{3. Modelo compartimental tipo DEB de Eriksen (2022)}
% --------------------------------------------------------------------

Este modelo es el que Eriksen (2024) denomina ``differential energy budget model'' (Kearney, 2021) adaptado a larvas de BSF. Divide la biomasa en tres compartimentos:

Este modelo divide la biomasa de la larva en tres compartimentos:

\begin{itemize}
    \item $A$ = reserva de asimilados (en estado cuasi-estacionario, análogo de $E$).
    \item $B$ = biomasa estructural (análogo de $V$).
    \item $L$ = lípidos de almacenamiento (reservas adicionales).
\end{itemize}

El sistema de ecuaciones diferenciales es:

\begin{align}
    \frac{dA}{dt} &= 0 = r_A - r_B - r_L - r_{CO_2,m} - r_{CO_2,B} - r_{CO_2,L}
    \label{eq:plos_A} \\[6pt]
    \frac{dB}{dt} &= r_B
    \label{eq:plos_B} \\[6pt]
    \frac{dL_{tot}}{dt} &= r_L + \delta_{L,B}\,r_B
    \label{eq:plos_L}
\end{align}

con las siguientes tasas constitutivas:

\begin{align}
    r_A &= a\,B
    \quad\text{(tasa de asimilación del alimento)} \label{eq:plos_rA} \\[4pt]
    r_{CO_2,m} &= m\,B
    \quad\text{(respiración de mantenimiento)} \label{eq:plos_rCO2m} \\[4pt]
    r_B &=
    \begin{cases}
        \mu_B\,B, & I \le 6 \\[-2pt]
        0, & I > 6
    \end{cases}
    \quad\text{(síntesis de biomasa estructural)} \label{eq:plos_rB} \\[4pt]
    r_{CO_2,B} &= Y_B\,r_B
    \quad\text{(costo respiratorio del crecimiento)} \label{eq:plos_rCO2B} \\[4pt]
    r_L &=
    \begin{cases}
        \dfrac{r_A - r_B - r_{CO_2,m} - r_{CO_2,B}}{1+Y_L}, & r_L \ge 0 \\[8pt]
        r_A - r_B - r_{CO_2,m} - r_{CO_2,B}, & r_L < 0
    \end{cases}
    \quad\text{(acumulación/movilización de lípidos)} \label{eq:plos_rL} \\[4pt]
    r_{CO_2,L} &=
    \begin{cases}
        Y_L\,r_L, & r_L \ge 0 \\[-2pt]
        0, & r_L < 0
    \end{cases}
    \quad\text{(costo respiratorio de lípidos)} \label{eq:plos_rCO2L}
\end{align}

El número de instar $I$ y la masa estructural máxima $B_{max}$ se predicen mediante ecuaciones logísticas de transición (Ec.\ 15--16 del artículo original). La tasa específica de crecimiento estructural $\mu_B$ se modela como:

\begin{equation}
    \mu_B = \frac{a-m}{1+Y_B}\left(1 - \frac{B}{B_{max}}\right)^\beta
    \label{eq:plos_muB}
\end{equation}

lo que introduce una reducción logística de la inversión en estructura conforme la larva se acerca a su masa estructural máxima $B_{max}$, forzando el excedente de asimilados hacia el almacenamiento de lípidos.

% --------------------------------------------------------------------
\section{4. Descripción de las variables y parámetros}
% --------------------------------------------------------------------

\begin{table}[htbp]
\centering
\caption{Variables de estado y sus interpretaciones}
\begin{tabular}{@{}cll@{}}
\toprule
\textbf{Símbolo} & \textbf{Significado} & \textbf{Equivalente DEB} \\
\midrule
$E$ & Energía en reserva & --- (estándar) \\
$V$ & Volumen estructural & --- (estándar) \\
$B_{dry}$ & Masa seca total por larva & $E+V$ (aprox.) \\
$B_{eff}$ & Asimilados no estructurales (reserva) & $E$ \\
$B_{str}$ & Masa estructural del cuerpo & $V$ \\
$A$ & Pool de asimilados (estado cuasi-estacionario) & $E$ \\
$B$ & Biomasa estructural & $V$ \\
$L_{tot}$ & Lípidos totales medidos (almacenamiento + estructurales) & reserva adicional \\
$TS$ & Suma de desarrollo (edad aparente) & --- \\
\bottomrule
\end{tabular}
\end{table}

\begin{table}[htbp]
\centering
\caption{Principales parámetros y tasas}
\begin{tabular}{@{}cll@{}}
\toprule
\textbf{Símbolo} & \textbf{Significado} & \textbf{Unidades típicas} \\
\midrule
$p_A$ & Tasa de asimilación & energía/time \\
$p_C$ & Tasa de utilización total de reserva & energía/time \\
$p_M=[p_M]V$ & Costo de mantenimiento & energía/time \\
$[E_G]$ & Costo de crecimiento por unidad de volumen & energía/volumen \\
$\kappa$ & Fracción de $p_C$ destinada a mantenimiento+crecimiento & adimensional \\
$\{p_{Am}\}$ & Tasa máxima de asimilación por unidad de superficie & energía/(area$\cdot$time) \\
$f$ & Función de alimentación (tipo Holling/Monod) & adimensional \\
$k_{inges}$ & Tasa específica de ingestión máxima & g alimento/(g larva$\cdot$s) \\
$\eta_{inges}$ & Eficiencia del alimento ingerido & adimensional \\
$k_{maint}$ & Tasa específica de mantenimiento y madurez & g/(g larva$\cdot$s) \\
$a$ / $a_{max}$ & Tasa específica de asimilación de alimento & d$^{-1}$ \\
$m$ & Tasa específica de producción de CO$_2$ por mantenimiento & g C/(g C$\cdot$d) \\
$Y_B$ / $Y_L$ & Costo de síntesis de biomasa/lípidos & adimensional \\
$\delta_{L,B}$ & Fracción de lípidos en la biomasa estructural & adimensional \\
\bottomrule
\end{tabular}
\end{table}

% --------------------------------------------------------------------
\section*{Observación final}
% --------------------------------------------------------------------

Ambos modelos empíricos para \textit{Hermetia illucens} conservan la filosofía del DEB (compartimentación de reservas y estructura, balance energético/másico, mantenimiento proporcional a la estructura), pero adaptan las ecuaciones a datos experimentales de masa seca y composición bioquímica. En particular, la ecuación (\ref{eq:pad_bdry}) de Padmanabha et al.\ y el sistema (\ref{eq:plos_A})--(\ref{eq:plos_L}) del artículo PLOS ONE son las formulaciones diferenciales principales extraídas directamente de los PDFs revisados.

\end{document}
